\section{The Abel--Jacobi map}
\label{mk-chap:AbelJacobi}

Let \(C\) be a smooth proper geometrically irreducible curve over a field \(k\),
let \(P \in C(k)\), and let \(\Jac(C)\) be the Albanese variety constructed in
Chapter~\ref{mk-chap:Jacobian}.  This chapter records the universal morphism and its
two defining properties.

\section{Definition}

\begin{definition}[Abel--Jacobi map]\label{mk-def:ofCurve}
  \lean{AlgebraicGeometry.Jacobian.ofCurve}
  \leanok
  \dcref{custom}
  \uses{mk-def:Jacobian, mk-def:IsAlbanese, mk-def:IsAlbanese_ofCurve,
    mk-thm:nonempty_jacobianWitness}

  The \emph{Abel--Jacobi map at \(P\)} is the universal pointed morphism
  \[
    \alpha_P \colon C \longrightarrow \Jac(C)
  \]
  supplied by the Albanese property of \(\Jac(C)\).
\end{definition}

\begin{remark}
[Picard description]
  \label{mk-rem:ofCurve_classical}
  \dcref{custom}
  Under the identification \(\Jac(C) = \Pic^0_{C/k}\), the morphism is
  \[
    \alpha_P(Q) = [\mathcal O_C(Q-P)].
  \]
  Equivalently, it is classified by
  \(\mathcal O_{C\times C}(\Delta-p_1^*P)\), regarded as a family of
  degree-zero line bundles over the second factor.
\end{remark}

\section{Pointed property}

\begin{lemma}[\(\alpha_P\) sends \(P\) to the identity]\label{mk-lem:comp_ofCurve}
  \lean{AlgebraicGeometry.Jacobian.comp_ofCurve}
  \leanok
  \dcref{custom}
  \uses{mk-def:ofCurve, mk-lem:IsAlbanese_comp_ofCurve,
    mk-thm:nonempty_jacobianWitness}

  The Abel--Jacobi map is pointed:
  \[
    P \circ \alpha_P = \eta_{\Jac(C)}.
  \]
\end{lemma}

\begin{proof}
  This is the pointed condition in the defining Albanese property of
  \(\Jac(C)\).  In the Picard description it is the equality
  \([\mathcal O_C(P-P)] = [\mathcal O_C]\).
\end{proof}

\begin{remark}
[Restriction of the universal line bundle]
  \label{mk-rem:comp_ofCurve_classical}
  \dcref{custom}
  Restricting \(\mathcal O_{C\times C}(\Delta-p_1^*P)\) along the section
  \((P,\mathrm{id}_C)\) gives the trivial line bundle on \(C\), which is the
  identity point of \(\Pic^0_{C/k}\).
\end{remark}

\section{Universal property}

\begin{theorem}[Albanese property of the Jacobian]\label{mk-thm:exists_unique_ofCurve_comp}
  \lean{AlgebraicGeometry.Jacobian.exists_unique_ofCurve_comp}
  \leanok
  \dcref{custom}
  \uses{mk-def:ofCurve, mk-def:IsAlbanese, mk-def:Jacobian,
    mk-lem:IsAlbanese_exists_unique_ofCurve_comp,
    mk-thm:nonempty_jacobianWitness}

  Let \(A\) be an abelian variety over \(k\).  For every morphism
  \(f \colon C \to A\) satisfying \(P \circ f = \eta_A\), there exists a
  unique morphism of \(k\)-schemes \(g \colon \Jac(C) \to A\) such that
  \[
    f = \alpha_P \circ g.
  \]
\end{theorem}

\begin{proof}
  By definition, \(\Jac(C)\) is an Albanese object for \((C,P)\), and
  \(\alpha_P\) is its universal pointed morphism.  Applying the universal
  factorisation property to \(f\) gives the required \(g\), while the
  uniqueness clause gives uniqueness.
\end{proof}
